\subsection{The output of the application, panel by panel}
Two layers show evapotranspiration and the crop requirement. \emph{Evaporation} follows the water from the pore to the atmosphere, and \emph{Crop water requirement} gives the net irrigation of each crop of the library on the soil of the explorer.

The picture on the left is the modelled soil in section, with its air-filled pores (white) and water-filled pores (blue); the graph shows the evaporation rate falling after wetting, from the first stage controlled by the atmosphere to the second controlled by the soil, and the cumulative loss.

\appfig{e1_0}{0.96}{Evaporation from the pore to the atmosphere: the modelled section of the soil, the evaporation rate and the cumulative loss through the drying spell.}

Six panels follow the flux in the soil, the suction, the water content and the pore-size distribution of the surface layer for the conventional setting and for the reduced and no-till settings.

\appfig{e1_1}{0.96}{Core-scale evaporation: from the matrix to the atmosphere, for the settings of the user.}

The sections give the water content and the direction of the flow on day 0.7; vapour leaves upward in the dry layer and liquid water moves up from below in the film of the pores.

\appfig{e1_2}{0.96}{Conventional (your setting): where the water moves, and the suction profile, day 0.7.}

The page begins with the crop and its coefficients and gives the net irrigation of the season for the soil and weather of the explorer, computed with $K_c\,ET_0$ and the trigger of the soil.

\appfig{crop_0}{0.96}{Crop water requirement: the crop, the net irrigation of the season and the daily need in the weather of the record.}

The depletion of the root zone is drawn against the readily available water; each irrigation returns it to field capacity, and the dates of the events are the schedule of the season.

\appfig{crop_1}{0.96}{Root-zone water and the irrigation schedule: depletion below field capacity, with irrigation and rain events of 2026.}

